\documentclass[10pt,twocolumn]{article}
\usepackage[margin=0.75in]{geometry}
\usepackage{amsmath,amssymb,booktabs,graphicx,microtype,url}
\usepackage[hidelinks]{hyperref}
\usepackage{xcolor}

\title{A Reproducibility-First Evaluation of Denoising-Time Skin-Tone
Steering in SDXL}
\author{Anonymous Author(s)\\\small Replace before submission}
\date{}

\newcommand{\method}{stepwise latent steering}

\begin{document}
\maketitle

\begin{abstract}
We study whether a linear direction estimated from paired portrait latents can
control visible skin tone in Stable Diffusion XL while preserving other image
properties. The proposed method computes a paired mean difference in the VAE
latent space, applies a spatial Gaussian mask, and distributes the resulting
perturbation across denoising steps. A single-seed pilot produces coherent
counterfactual images and motivates a confirmatory evaluation protocol.
However, the checked-in pilot lacks valid face-identity and pose measurements,
and its historical ``background SSIM'' values reduce to whole-image SSIM when
face detection is unavailable. We therefore present the result as feasibility
evidence, not evidence of identity preservation, disentanglement, or bias
mitigation. Alongside the method, we contribute a claim--evidence register,
deterministic encoding and cache behavior, provenance capture, missing-metric
failure rules, a frozen multi-seed comparison against prompt-only, post-hoc,
and unmasked baselines, and a fail-closed precollection readiness gate. The
work explicitly treats skin tone as a visual attribute and not as a proxy for
race or ethnicity. Confirmatory collection remains blocked pending external
target-metric validation, supported runtime and artifact verification, exact
generation provenance, and explicit approval.
\end{abstract}

\section{Introduction}

Latent diffusion models synthesize high-resolution images by denoising in a
compressed representation rather than directly in pixel space
\cite{rombach2022ldm}. SDXL extends this family with a larger U-Net and richer
conditioning \cite{podell2023sdxl}. These models offer expressive control but
also make it difficult to isolate one visual attribute without changing
identity, pose, or composition.

Linear semantic directions have been effective in GAN latent spaces
\cite{shen2020interfacegan}. We ask whether an analogous direction, estimated
from paired SDXL VAE encodings, can be used during diffusion sampling. Our
technical hypothesis is that a denoiser can absorb small perturbations injected
throughout sampling more coherently than a single post-hoc edit. Our scientific
hypothesis is narrower: at a matched change in measured skin tone, stepwise
masked injection will preserve non-target properties better than relevant
baselines.

We make three contributions. First, we specify the paired estimator and
denoising-time intervention. Second, we audit a working pilot and distinguish
observed evidence from unsupported claims. Third, we provide a confirmatory
protocol and research artifact designed to report provenance, uncertainty,
missingness, and failures.

\section{Construct and scope}

Race and ethnicity are social identities, not one-dimensional visual
attributes. The intervention in this study is therefore called a
\emph{skin-tone direction}. It must not be interpreted as changing, estimating,
or classifying race. Even datasets designed to balance face categories expose
the difficulty and consequences of reducing social categories to visual labels
\cite{karkkainen2021fairface}. The present method is intended for controlled
generative-model research and auditing, not profiling or decisions about
people.

\section{Method}

Let $E(x)$ be the deterministic posterior-mode encoding of an image under the
SDXL VAE. Given $n$ paired portraits $(x_i^{L},x_i^{D})$ generated with matched
seed and composition but different skin-tone descriptors, the raw direction is
\begin{equation}
  v = \frac{1}{n}\sum_{i=1}^{n}\left[E(x_i^{D})-E(x_i^{L})\right].
\end{equation}
The pairing attempts to cancel seed-specific facial structure, although prompt
conditioning and illumination remain potential confounders.

We multiply $v$ by a spatial mask $M\in[0,1]^{H\times W}$ centered on the face:
\begin{equation}
  M(p)=w_e+(w_c-w_e)\exp\left[-3\left(\frac{\lVert p\rVert_2}{r}\right)^2\right],
  \qquad \tilde v=M\odot v,
\end{equation}
where the pilot uses center weight $w_c=1$, edge weight $w_e=0.3$, and radius
$r=1$. For $T$ reverse-diffusion steps and steering magnitude $\alpha$, the
pipeline updates the scheduler latent after each step by
\begin{equation}
  z_t \leftarrow z_t + \frac{\alpha}{T}\tilde v.
\end{equation}
The implementation uses 25 DPM-Solver++ multistep updates, motivated by fast
diffusion ODE solvers \cite{lu2022dpmsolver}, with classifier-free guidance 7.5.

\section{Evaluation design}

The target response and preservation outcomes must be measured separately.
The primary planned outcome is face-embedding similarity, following the general
embedding approach of FaceNet \cite{schroff2015facenet}, at a matched continuous
skin-tone change. Secondary outcomes are LPIPS perceptual distance
\cite{zhang2018lpips}, background-only SSIM \cite{wang2004ssim}, head-pose
drift, face-detection failure, and monotonicity.

The confirmatory design uses held-out seeds and four methods: prompt-only
generation, post-hoc latent addition, unmasked stepwise injection, and masked
stepwise injection. The planned configuration contains 30 held-out generation
seeds and five alpha values. Comparisons are paired within seed. We will report
means, standard deviations, medians, detector failures, and 95\% seed-level
bootstrap confidence intervals; secondary confirmatory tests use Holm
correction. A run cannot pass when a required metric is missing.

Before confirmatory collection, all seven criteria in
\texttt{tmlr\_collection\_readiness\_v1} must pass: held-out calibrated-
instrument validation of the target metric, synthetic illumination and color
sensitivity checks, the supported Linux metric runtime, checksum-verified
metric weights, immutable SDXL and direction provenance, the storage/privacy
policy, and an approval record bound to those exact inputs. The current
non-generative validation passes only the deterministic synthetic checks and
the versioned storage/privacy policy. It therefore provides implementation
evidence but no confirmatory outcome or external-validity claim.

\section{Pilot result}

Figure~\ref{fig:pilot} shows one synthetic base portrait and four steered
outputs. The current metadata records LPIPS values between 0.134 and 0.280 and
historical SSIM values between 0.798 and 0.864 (Table~\ref{tab:pilot}). No valid
face-embedding, landmark, or pose values were recorded. Because the historical
background routine fell back to an all-background mask after face-detection
failure, these SSIM values characterize the whole masked canvas rather than a
validated background region. They are retained for auditability but are not
treated as evidence of background preservation.

\begin{figure*}[t]
  \centering
  \includegraphics[width=0.95\textwidth]{../experiments/results/final_grid.png}
  \caption{Single-seed pilot. Labels embedded in this historical figure should
  be interpreted with the corrected evidence record: identity and pose were not
  available, and the displayed ``background SSIM'' is not a validated
  background-only measurement.}
  \label{fig:pilot}
\end{figure*}

\begin{table}[t]
\centering
\caption{Checked-in pilot metadata. Dashes denote unavailable measurements.}
\label{tab:pilot}
\begin{tabular}{@{}rrrr@{}}
\toprule
$\alpha$ & LPIPS $\downarrow$ & historical SSIM $\uparrow$ & face sim. \\
\midrule
$-1.50$ & 0.280 & 0.821 & -- \\
$-0.75$ & 0.242 & 0.847 & -- \\
$+0.75$ & 0.134 & 0.864 & -- \\
$+1.50$ & 0.204 & 0.798 & -- \\
\bottomrule
\end{tabular}
\end{table}

\section{Limitations}

The pilot has one evaluation seed, six training pairs, no held-out target
measurement, and no valid biometric or pose metrics. Synthetic pairs may encode
prompt wording, brightness, age, gender presentation, hair, or facial features
alongside skin tone. A center Gaussian is not a semantic skin mask. The same
generator creates training and evaluation images, so model-specific artifacts
may dominate. Face embeddings and skin-tone estimators can themselves exhibit
demographic error. The relative-cheek CIELAB measure reduces a uniform exposure
shift in a synthetic fixture but remains confounded by local face illumination;
it has not been validated against calibrated measurements in the intended
portrait domain. Finally, the additive update is scheduler-dependent and has
not been validated across model or library revisions.

\section{Conclusion}

The pilot establishes that denoising-time latent injection is implementable and
can yield a coherent qualitative sweep. It does not establish identity
preservation, disentanglement, or bias mitigation. The released research
scaffold turns those assertions into falsifiable, prespecified questions and
records the evidence needed to answer them.

\bibliographystyle{plain}
\bibliography{references}
\end{document}
